\section{Discussion}
\label{sec:discussion}

\subsection{The dissociation, and the production stance it licenses}

The result we consider most consequential for practice is the SECOM row of Table~\ref{tab:t3}.
A soft sensor that cannot be made accurate is a common industrial outcome: the
informative variables are not instrumented, the process is not excited, the physics is
not available. The usual response is either to abandon the sensor or to operate it
with confidence statistics inherited from a validation set it no longer resembles. The
framework offers a third position: deploy the calibration machinery even when the
point model fails, and let the intervals tell the truth. On SECOM the corrected + ACI
construction holds 0.910/0.915 coverage at 37\% narrower width than the over-covering
static baseline \emph{around a point model with negative $R^2$}. An operator reading those
intervals learns, correctly, that the sensor knows little, which is actionable in a
way that a confidently wrong point estimate is not. Conformal validity is
model-agnostic; the decision to trust a soft sensor should be made on its interval
width, not on the existence of a prediction.

\subsection{Where the performance comes from}

Holding the pipeline fixed across three datasets decomposes the credit. Point accuracy
comes from the physics anchors: the same protocol that yields test $R^2$ of 0.48--0.86 on
the debutanizer and working sensors on TEP produces --1.84 on SECOM, where no
thermodynamic structure exists to encode (Sections 5.1, 5.5). Robustness, the
difference between a usable and a lottery-dependent sensor \emph{within} a dataset, comes
from the validation discipline: blocked CV exposes the regime spread that adjacent
validation averages away, and the one-SE rule converts that exposure into parsimonious
selection (the k = 126 kitchen sink and the SECOM path explosion are the same failure
at two scales). Honesty comes from the conformal layer, and only from it: neither the
physics nor the bias update carries a coverage guarantee. The components are
complementary, and the composition order matters: bias correction before conformal
calibration (so the intervals price residual uncertainty, not removable offset), and
drift alarming on the corrected residual (so alarms mean ``intervention needed'', not
``the EWMA is working'').

\subsection{The linear scope is a design choice, and what it costs}

Every point model in this paper is linear. An obvious objection is whether a
gradient-boosted or deep model would rescue SECOM. We answer in two parts. First, the
negative-control logic \emph{requires} holding the model class fixed: the claim ``physics
anchors are the accuracy difference'' is identified only because the estimator is
identical across datasets; swapping in a nonlinear learner on the hard dataset would
confound model capacity with feature information. Second, the linear choice is what
makes the rest of the stack cheap and auditable: the bias update and conformal
bookkeeping are O(1) per label, the per-estimate computation is microseconds, and Luo's
degeneracy result (Section 5.4) is itself a consequence of source linearity that the
literature obscures. The cost is stated plainly: the SECOM sensor is not deployable
for point prediction, and we do not know how much of its failure a nonlinear model
would recover. That experiment, which also re-opens Luo's method since its collapse
is specific to linear sources, is the designated extension, and SECOM now provides
the exact dataset on which to run it.

\subsection{Limitations}

Stated as facts, not hedges. (i) The migration result is within-process regime
migration; cross-process parameter migration is inapplicable by construction
(disjoint input spaces), and what transfers across processes is the methodology.
(ii) The TEP regimes are feed-ratio operating points generated under deliberate
excitation, not the canonical product modes; this is sufficient for the within-process claim,
not a canonical-benchmark claim. (iii) The fail-enrichment analysis on SECOM is
underpowered (9 failing lots in the test window) and is reported as inconclusive.
(iv) The OSBC data-efficiency figure is sensitive to evaluation conditions (pool
subsampling, the presence of the online layer): it measured \textasciitilde3.3$\times$ in early runs and
1.0$\times$/4.0$\times$ under the frozen protocol, which is one reason the paper commits to provenance-
stamped evidence and reports the frozen values. (v) The online implementation runs as a single instance with persisted calibration state; coordinating several parallel instances would require sharing that state, which is documented but not built. (vi) The label delay is treated as fixed and
known; random or drifting delays are handled in the literature (Section 2.1) and
would compose, but are not evaluated here. (vii) The virtual-metrology target on SECOM
is our construction; the selection criteria and audit table are published so
that the choice can be attacked.

\subsection{Practitioner's recipe and future work}

For a practitioner the framework reduces to an ordering: diagnose the transport lag
before judging any model; encode what physics is available as features and let k stay
small; select with blocked CV and the one-SE rule, treating negative folds as signal;
correct drift with the open-loop EWMA at the documented delay; calibrate with ACI on
the corrected residuals and alarm on them too; and read interval width as the trust
metric. Each step is justified by a failure it removed in Section 5.

Beyond the nonlinear extension (\S 6.3), three directions follow directly. Dynamic
delay estimation \citep{wang2021dtde} is a drop-in upgrade where analyzer delays drift.
The online implementation was designed for connection to plant data infrastructure (OPC~UA/MQTT ingestion was prototyped and then gated out of scope); closing that loop turns the reference implementation into a plant pilot. And the evidence-freeze
discipline used to produce this paper, in which evaluation scripts emit provenance-stamped
artifacts that figures and claims consume, proved valuable enough during writing
(it caught two would-be mischaracterizations before submission) that we would propose
it as practice independent of the framework itself.
